% A1-anexo-variantes.tex — Anexo: las dos variantes de estrés de memoria de Q2.
% Requisito R-16.1: cobertura parcial, decidida y justificada, FUERA del CSV de
% 36 celdas, que no se toca.

\section{Anexo: las dos variantes de estrés de memoria de \texttt{Q2}}
\label{sec:variantes-anexo}

\subsection{Por qué este anexo existe, y por qué es parcial}

El manual (§5.2) pide ejecutar \texttt{Q2} con dos configuraciones de memoria. El
conteo de escenarios que fija el propio manual son 36 celdas, y con 36 el esquema
del CSV \textbf{no tiene ninguna columna donde distinguir la variante}: añadir una
está prohibido explícitamente. Las dos exigencias no caben juntas.

La decisión tomada fue medir la segunda variante \textbf{en anexo}, fuera del CSV
de las 36 celdas, y solo en \textbf{\texttt{PROFILE\_C\_SEVERE}}: los dos motores y
las tres concurrencias, 6 celdas y 18 repeticiones. El razonamiento:

\begin{itemize}
  \item \texttt{PROFILE\_C} es el perfil con \textbf{menos memoria} (1~GiB). Si las
        dos variantes no se diferencian ahí, no se diferencian en ninguno.
  \item Por el plan ya se sabía que el ordenamiento de \texttt{Q2} ocupa
        \textbf{273~kB} (sección~\ref{sec:planes}), muy por debajo tanto de los
        4~MB de la variante restrictiva como de los 64~MB de la permisiva.
\end{itemize}

\textbf{La cobertura es parcial y se declara como tal}: no se midió
\texttt{Q2\_MEM\_HIGH} en los perfiles A y B. Los datos están en
\texttt{results/f8\_anexo/measurements.csv} (18 filas, validadas con el mismo
esquema de 14 columnas).

\subsection{Resultado: las dos variantes no se diferencian}

\begin{table}[htbp]
  \centering
  \caption{\texttt{Q2} en \texttt{PROFILE\_C\_SEVERE}: variante restrictiva
           (\texttt{work\_mem} 4~MB / \texttt{allowDiskUse} falso) frente a
           permisiva (64~MB / verdadero). Throughput medio de 3 repeticiones.}
  \label{tab:variantes}
  \small
  \begin{tabular}{lrrrrl}
    \toprule
    Motor & c & Restrictiva & Permisiva & Diferencia & Volcado \\
    \midrule
    MongoDB    & 1  & 0.0533 & 0.0444 & $-16.7$~\% & ninguno en ninguna \\
    MongoDB    & 10 & 0.1000 & 0.1000 & $0.0$~\%   & ninguno en ninguna \\
    MongoDB    & 50 & 0.1667 & 0.1667 & $0.0$~\%   & ninguno en ninguna \\
    PostgreSQL & 1  & 0.5722 & 0.6622 & $+15.7$~\% & ninguno en ninguna \\
    PostgreSQL & 10 & 0.6000 & 0.6222 & $+3.7$~\%  & ninguno en ninguna \\
    PostgreSQL & 50 & 0.5611 & 0.5744 & $+2.4$~\%  & ninguno en ninguna \\
    \bottomrule
  \end{tabular}
\end{table}

Las dos diferencias mayores (un $-16.7$~\% en MongoDB y un $+15.7$~\% en
PostgreSQL, las dos con un cliente) \textbf{apuntan en direcciones opuestas}, que
es la forma que tiene el ruido, no un efecto del parámetro. Las otras cuatro están
entre el 0 y el 3.7~\%, por debajo de la dispersión propia de esas celdas. Las 18
repeticiones salieron \texttt{ok} y \textbf{ninguna volcó a disco}, ni con
\texttt{allowDiskUse: true} en MongoDB ni con 4~MB de \texttt{work\_mem} en
PostgreSQL.

\subsection{Conclusión}

\textbf{La premisa del manual no se cumple.} Se esperaba que la variante
restrictiva forzara volcado a disco y que las dos configuraciones dieran resultados
distintos; medidas las dos en el perfil más exigente, no se diferencian y ninguna
vuelca. La causa está en el plan y no en el parámetro: el \texttt{LIMIT 1000}
acota el ordenamiento a 273~kB, así que el límite de memoria que el manual varía
nunca se alcanza. Se reporta como hallazgo, no como fallo del experimento.
